\uuid{DzLc}
\chapitre{Statistique}
\niveau{L2}
\module{Analyse de données}
\sousChapitre{Tests d'hypothèses, intervalle de confiance}
\titre{Test d'indépendance}
\theme{tests d'hypothèses}
\auteur{}
\datecreate{2022-10-19}
\organisation{AMSCC}
\difficulte{}
\contenu{
\texte{Les quantiles sont à gauche : $q_{\nu;\gamma}$ pour la loi $\chi^2(\nu)$, avec une probabilité $\gamma$ à gauche du quantile.}
\texte{Une enquête historique auprès de 2089 personnes croise l'âge et la réponse à une affirmation sur le sida. On suppose les répondants indépendants et les conditions d'échantillonnage adaptées au test de Pearson.
\begin{center}\begin{tabular}{lrrrr}\hline Âge & D'accord & Indécis & Pas d'accord & Total \\ \hline
$[18;30[$ & 43 & 116 & 365 & 524 \\
$[30;40[$ & 36 & 116 & 273 & 425 \\
$[40;50[$ & 32 & 95 & 217 & 344 \\
$[50;60[$ & 38 & 114 & 167 & 319 \\
60 et plus & 67 & 160 & 250 & 477 \\
Total & 216 & 601 & 1272 & 2089 \\ \hline\end{tabular}\end{center}}
\begin{enumerate}
\item \question{Tester l'indépendance de l'âge et de la réponse au seuil de $1\%$.}
\indication{Sous $H_0$, calculer $e_{ij}=\frac{o_{i\cdot}o_{\cdot j}}{2089}$.}
\reponse{Les effectifs attendus sont :
\begin{center}\begin{tabular}{lrrr}\hline Âge & D'accord & Indécis & Pas d'accord \\ \hline
$[18;30[$ & 54,18095 & 150,75347 & 319,06558 \\
$[30;40[$ & 43,94447 & 122,27142 & 258,78411 \\
$[40;50[$ & 35,56917 & 98,96793 & 209,46290 \\
$[50;60[$ & 32,98420 & 91,77549 & 194,24031 \\
60 et plus & 49,32121 & 137,23169 & 290,44710 \\ \hline\end{tabular}\end{center}
Ils sont tous supérieurs à 5. On trouve $q_{\mathrm{obs}}=\sum\frac{(o_{ij}-e_{ij})^2}{e_{ij}}\simeq45{,}97114$ avec $(5-1)(3-1)=8$ ddl.
Le seuil $q_{8;0{,}99}\simeq20{,}090$ est dépassé et $p_{\mathrm{val}}\simeq2{,}41\times10^{-7}$ : on rejette l'indépendance. Une association est mise en évidence, sans conclusion causale.}
\item \question{Au seuil de $5\%$, comparer les deux dernières classes d'âge, puis les moins de 30 ans aux 60 ans et plus.}
\reponse{Chaque comparaison demande de recalculer les marges sur les seules lignes retenues.
Pour les deux dernières classes, les attendus sont $(42{,}07915;109{,}80653;167{,}11432)$ et $(62{,}92085;164{,}19347;249{,}88568)$.
Avec 2 ddl, $q_{\mathrm{obs}}\simeq0{,}927260$ et $p_{\mathrm{val}}\simeq0{,}628996$ : non-rejet. On ne met pas en évidence une différence de distributions ; on n'a pas prouvé leur égalité.
Pour les extrêmes, les attendus sont $(57{,}58242;144{,}47952;321{,}93806)$ et $(52{,}41758;131{,}52048;293{,}06194)$.
On obtient $q_{\mathrm{obs}}\simeq31{,}61783$, 2 ddl et $p_{\mathrm{val}}\simeq1{,}36\times10^{-7}$ : rejet.
Ces comparaisons sont des analyses distinctes ; une exploration de nombreuses paires demanderait de tenir compte de la multiplicité des tests.}
\end{enumerate}
}
